

\begin{Exercice}[10 minutes]\textbf{Sean}

Dans cet exercice nous aimerions que le processeur calcule et enregistre le produit de 3 fois $10_{10}$ ($00001010_2$) dans le registre 11. Le program counter (PC) commence à \textbf{00000100} et les registres sont initialisés comme suit:

\begin{center}
\begin{tabular}{|c|l|}
\hline
\textbf{Registre} & \textbf{Valeur} \\
\hline
00 & 00001010 \\ \hline
01 & 00000000 \\ \hline
10 & 00000000 \\ \hline
11 & 00000000 \\ 
\hline
\end{tabular}
\end{center}

Les opérations disponibles sont:

\begin{center}
\begin{tabular}{|c|c|}
\hline
\textbf{Numéro} & \textbf{Instruction} \\
\hline
00 & MOV \\ \hline
01 & XOR \\ \hline
10 & ADD \\ \hline
11 & SUB \\ 
\hline
\end{tabular}
\end{center}

\bigskip % Creates vertical space, equivalent to Typst's \ \

% --- Part a ---
\noindent % Prevents indentation
\textbf{a.} Complétez le tableau ci-dessous avec les instructions en binaire nécessaires pour effectuer cette tâche.

\begin{center}
\begin{tabular}{|l|l|}
\hline
\textbf{Adresse} & \textbf{Valeur} \\
\hline
00000100 & \\ \hline
00000101 & \\ 
\hline
\end{tabular}
\end{center}

\bigskip

% --- Part b ---
\noindent % Prevents indentation
\textbf{b.} Donnez le contenu des registres après l'exécution de la première instruction et après l'exécution de la deuxième instruction.

\begin{center}
\begin{tabular}{|c|l|l|}
\hline
\textbf{Registre} & \textbf{Après 1ère instruction} & \textbf{Après 2ème instruction} \\
\hline
00 & & \\ \hline
01 & & \\ \hline
10 & & \\ \hline
11 & & \\ 
\hline
\end{tabular}
\end{center}

\begin{solution}
\textbf{a.}
\begin{center}
\begin{tabular}{|l|l|}
\hline
\textbf{Adresse} & \textbf{Valeur} \\
\hline
00000100 & 10100000\\ \hline
00000101 & 10110010\\ 
\hline
\end{tabular}
\end{center}

\bigskip

\textbf{b.}

\begin{center}
\begin{tabular}{|c|l|l|}
\hline
\textbf{Registre} & \textbf{Après 1ère instruction} & \textbf{Après 2ème instruction} \\
\hline
00 & 00001010 & 00001010 \\ \hline
01 & 00000000 & 00000000 \\ \hline
10 & 00010100 & 00010100 \\ \hline
11 & 00000000 & 00011110 \\ 
\hline
\end{tabular}
\end{center}

\end{solution}
\end{Exercice}

\begin{Exercice}[15 minutes]\textbf{Sean}
Quel est le résultat de la soustraction hexadécimale suivante en base 16: 694F - 5A3C?

\begin{solution}
0F13
\end{solution}
\end{Exercice}


\begin{Exercice}[10 minutes]\textbf{Aurore}
Selon le modèle de Von Neumann, partant d'un program counter valant 101010 et d'une mémoire contenant les instructions suivantes:
\begin{table}[h]
\centering
\begin{tabular}{|c|c|}
\hline
\textbf{Adresse} & \textbf{Instruction} \\
\hline
101001 & 01000111 \\ \hline
101010 & 01111000 \\ \hline
101011 & 10001101 \\ \hline
101100 & 10010110 \\ 
\hline
\end{tabular}
\end{table}

Dans quel registre le résultat de la prochaine opération va-t-il être stocké?

\mcqlist{00, 01, 10, 11}

\begin{solution}
11
\end{solution}

\end{Exercice}

\begin{Exercice}[10 minutes]\textbf{Aurore} Quel est le résultat de l'opération suivante : $395_{16}$ + $1011101_{2}$?
\mcqlist{$1011_{10}$, $3F2_{16}$, $0011 1110 0011_2$, $1010_2$}

\begin{solution}
$3F2_{16}$
\end{solution}
\end{Exercice}

\begin{Exercice}[10 minutes]\textbf{Maxim} Quelle est la valeur en base 10 de l’opération suivante : $128_9$ - $245_6$?
\mcqlist{9,6,0,3}
\begin{solution}
6
\end{solution}
\end{Exercice}

\begin{Exercice}[10 minutes]\textbf{Maxim}
Convertir le nombre $127_{10}$ en base 8.
\mcqlist{87,177,71,180}
\begin{solution}
$177_8$
\end{solution}
\end{Exercice}

\begin{Exercice}[10 minutes]\textbf{Léa}
Quel est le résultat de l'opération suivante $67_8$ + $110111_2$ en base 10? 
\mcqlist{110, 78, 122, 106}
\begin{solution}
    110
\end{solution}
\end{Exercice}

\begin{Exercice}[10 minutes]\textbf{Léa}
En suivant le modèle de von Neumann, l'Instruction Register (IR) a une valeur de 11000111, les registres contiennent les valeurs suivantes:
\begin{table}[h!]
\centering
\begin{tabular}{|c|c|}
\hline
\textbf{Registre} & \textbf{Valeur} \\ \hline
00 & 00000000 \\  \hline
01 & 00000101 \\  \hline
10 & 00000110 \\  \hline
11 & 00000010 \\ \hline
\end{tabular}
\end{table}

\begin{center}
\begin{tabular}{|c|c|}
\hline
\textbf{Numéro} & \textbf{Valeur} \\ \hline
00 & MOV \\  \hline
01 & XOR \\  \hline
10 & ADD \\  \hline
11 & SUB \\ \hline
\end{tabular}
\end{center}
Quelle est la valeur du résultat de l'opération?
\mcqlist{$00000111_2$, $00001011_2$, $00000011_2$, $00000001_2$}
\begin{solution}
$00000011_2$ est stocké dans le registre 00.
\end{solution}
\end{Exercice}

% \begin{Exercice}[10 minutes]\textbf{Léa}
% Quelle est la valeur du résultat de l'opération?
% \mcqlist{$111_2$, $1011_2$, $011_2$, $0001_2$}
% \end{Exercice}

\begin{Exercice}[10 minutes]\textbf{Reda}
Que représentent les deux premiers bits (partant de gauche) de l'Instruction Register? (exemple: 11001001)?
\mcqlist{L'adresse du registre dans lequel il faudra mettre le résultat,
L'identificateur de l'opération à exécuter,
L'adresse du registre dans lequel est contenue la valeur sur laquelle faire l'opération,
Ils n'ont aucune signification}

\begin{solution}
Option B.
\end{solution}
\end{Exercice}

\begin{Exercice}[10 minutes]\textbf{Reda}
Quel est le résultat de la soustraction $11011001_2$ - $00110110_2$ en base 10?
\begin{solution}
163
\end{solution}
\end{Exercice}

\begin{Exercice}[10 minutes]\textbf{Amina}
Quel est le résultat de l'opération suivante : $1011010_2$ - $1000111_2$?
\mcqlist{$0001111_2$, $0010011_2$, $0100011_2$, $0010101_2$}
\begin{solution}
$0010011_2$
\end{solution}
\end{Exercice}


\begin{Exercice}[10 minutes]\textbf{Amina}
Quel est le résultat en base 10 de l'opération suivante : $47_{10}$ + $DEF_{16}$ ?
\mcqlist{3614, 3478, 3589, 3651}
\begin{solution}
    3614
\end{solution}
\end{Exercice}


\begin{Exercice}[10 minutes]\textbf{Léa}
Écrivez le complément à 2 de $21_{10}$ sur 8 bits:
\mcqlist{1110 1011, 1110 1010, 0001 0101, 0001 0100}

\begin{solution}
Option A. 
\end{solution}
\end{Exercice}


\begin{Exercice}[10 minutes]\textbf{Amina}
Quelle est la représentation du complément à 2 (two complement) de -143 sur 8 bits?

\mcqlist{10010001, 01110001, 01100001, Aucune des réponses n'est correcte}

\begin{solution}
La valeur la plus négative sur 8 bits en complément à deux est $-2^7=-128$. Aucune des réponses n'est correcte.
\end{solution}
\end{Exercice}
